\documentclass[12pt, letterpaper]{article}
\usepackage[utf8]{inputenc}
\usepackage[T1]{fontenc}
\usepackage{lmodern}
\usepackage[
mathjax,
ImagesDirectory=imgs,
ImagesName={lateximg-}
]{lwarp}

\usepackage[margin=1in]{geometry}
\usepackage{setspace}
\doublespacing
\usepackage{csquotes}
\usepackage{textalpha}
\usepackage{alphabeta}

\setlength{\marginparwidth}{2cm}
\usepackage[disable]{todonotes}

\DeclareUnicodeCharacter{2227}{\ensuremath{\wedge}}
\DeclareUnicodeCharacter{00AC}{\ensuremath{\neg}}
\DeclareUnicodeCharacter{2192}{\ensuremath{\rightarrow}}

\usepackage{amsmath, amssymb, graphicx}
\usepackage{textcomp}
\usepackage{xcolor}
\usepackage{tikz}
\usepackage{float}
\usetikzlibrary{calc, positioning}
\usetikzlibrary{decorations.pathreplacing}
\usepackage[font=small]{caption}
\usepackage{mathtools}
\usepackage{comment}

\usepackage[notes,backend=biber]{biblatex-chicago} 
\addbibresource{references.bib}
\usepackage[colorlinks=true, linkcolor=blue, citecolor=blue, urlcolor=blue]{hyperref}

\widowpenalty=10000
\clubpenalty=10000
\DeclareRobustCommand{\code}[1]{%
	\begingroup
	\setlength{\fboxsep}{2pt}%
	\colorbox{gray!15}{\ttfamily\small #1}%
	\endgroup
}
\newcommand{\supcom}{%
	\textsuperscript{,}
}
\newcommand{\multicomment}[1]{}
\usepackage{listings}

\newcommand{\customitem}[2]{%
	\noindent\textbf{#1}\\[0.5ex]%
	#2\vspace{1ex}%
}

\lstnewenvironment{lispcode}
{\lstset{
		language=Lisp,
		alsoletter={-,:},
		morekeywords={
			defparameter, defun, loop, for, from, to, do,
			multiple-value-bind, if, when, format, sleep, ql:quickload,
			make-atomic-ref, atomic-ref-value, cas, atomic-incf, memory-barrier,
			make-queue, queue-push, queue-pop, queue-empty-p,
			make-bounded-queue, bounded-queue-capacity, bounded-queue-push, 
			bounded-queue-pop, bounded-queue-empty-p, bounded-queue-full-p, 
			bounded-queue-push-batch, bounded-queue-pop-batch,
			make-spsc-queue, spsc-push, spsc-pop
		},
		basicstyle=\ttfamily\small\linespread{1}\selectfont\raggedright,
		columns=fullflexible,
		keepspaces=true,
		upquote=true,
		keywordstyle=\color{blue}\bfseries,
		commentstyle=\color{green!60!black}\itshape,
		stringstyle=\color{red},
		breaklines=true,
		showstringspaces=false,
		aboveskip=1.5em,
		belowskip=1.5em,
		xleftmargin=0pt,
		backgroundcolor=\color{gray!10},
		frame=single,
		framerule=1.5pt,
		rulecolor=\color{black}
}}
{}

\lstnewenvironment{bashcode}
{\lstset{
		language=bash,
		alsoletter={-},
		morekeywords={
			sudo, git, apt, apt-get, pacman, echo, cd, ls, mkdir, 
			rm, mv, cp, export, source, make, cmake, ql
		},
		basicstyle=\ttfamily\small\linespread{1}\selectfont\raggedright,
		columns=fullflexible,
		keepspaces=true,
		upquote=true,
		keywordstyle=\color{blue}\bfseries,
		commentstyle=\color{green!60!black}\itshape,
		stringstyle=\color{red},
		breaklines=true,
		showstringspaces=false,
		aboveskip=1.5em,
		belowskip=1.5em,
		xleftmargin=0pt,
		backgroundcolor=\color{gray!10},
		frame=single,
		framerule=1.5pt,
		rulecolor=\color{black}
}}
{}
% n.b, I had to run sudo pacman -S texlive-plaingeneric ...
% extra/texlive-plaingeneric 2026.1-1 (texlive)
% usr/share/texmf-dist/tex/latex/tracklang/tracklang.sty
% datetime2 requires tracklang.sty ...
% Verify it can be found with kpsewhich tracklang.sty
\usepackage[style=iso]{datetime2}

\newcommand{\lastedited}{%
	Last edited: \DTMnow\ UTC%
}
\date{\lastedited}
\setcounter{FileDepth}{1}
\title{cl-freelock Manual}
\author{Andrew D. France}
\date{\lastedited}

\begin{document}
	
	\maketitle
	\newpage
	
	\tableofcontents
	\newpage
	\section{Tutorial}
	\subsection{Getting Started: Installation}
	Once available on quicklisp, you will be able to simply run:
\begin{lispcode}
(ql:quickload :cl-freelock)
\end{lispcode}
	For now, clone the repository into your Quicklisp local-projects directory:
\begin{bashcode}
cd ~/quicklisp/local-projects/
git clone https://github.com/ItsMeForLua/cl-freelock.git
\end{bashcode}
	Then, load the system in your REPL:
\begin{lispcode}
(ql:quickload :cl-freelock)
\end{lispcode}

	\subsection{Basic Usage}
	In this section, we will go over the basic usage of unbounded Multi-Producer Multi-Consumer (MPMC), bounded MPMC, bounded Single Producer Single Consumer (SPSC), and error handling.
	
	\textbf{Additional Information}
	\begin{enumerate}
		\item \texttt{cl-freelock} allows the use of \code{fl} when calling its functions.
		\subitem  For example: \code{(fl:queue-push *work-queue* i)}
		\item The API reference for cl-freelock is mostly written with SBCL in-mind. However, the differences between implementations regarding the code presented in this documentation will be quite minor.
	\end{enumerate}
	
	
	
	\subsubsection{MPMC Unbounded Queue (Michael-Scott)}
	The algorithm used for unbounded queues is based off of the Michael-Scott algorithm. We specifically use Compare-And-Swap (CAS) operations to link nodes and swing the head and tail pointers.
	
	We identify three main reasons one would use this particular queue type:
	\begin{enumerate}
		\item When the workload is highly variable in predictability, such that you cannot guarantee a maximum capacity for your queue ahead of time.
		\item You have a thread pool architecture with multiple producer threads that are throwing jobs onto the queue, and multiple worker threads are simultaneously pulling jobs off. i.e., An N-to-N architecture.
		\item If your project requires immunity from deadlocks.
	\end{enumerate}
	\textbf{Producers}
	A producer thread works by calling \code{(queue-push *work-queue* item)}. Because the producer thread is lock-free, the procedure will never block a thread. An example of defining a producer:
\begin{lispcode}
(defun producer ()
(loop for i from 1 to 1000 do
(fl:queue-push *work-queue* i)
(sleep 0.001)))
\end{lispcode}
	n.b., It is very important you consider the reasoning for adding a sleep in this example. The sleep is just to simulate the time it may take to fetch the next item. It is not needed in the procedure within normal use-cases; however, if you simply ran the example in a fresh REPL, you would risk locking up your entire environment.\footnote{This is called ``CPU Starvation'', which is when a procedure consumes 100\% of a systems CPU core.}\supcom\footnote{Concurrent programming is a matter of managing the interleaving of different processes as they interact with a shared state --- hence, if we ran a loop from i to 1000 without sleep, most modern CPUs would actually be so fast that the producer may finish pushing all 1000 items before a consumer thread has a chance to wakeup.}
	
	\textbf{Consumers}
	A consumer thread works by calling \code{(queue-pop *work-queue*)}. This procedure does not return lisp conditions; instead, \code{queue-pop} returns two specific values:
	\begin{enumerate}
		\item The \code{item} itself or \code{NIL} if empty.
		\item A \code{success} boolean, \code{T} if an item was grabbed, \code{NIL} if the queue was empty.
	\end{enumerate}
	n.b., If a queue happens to be empty, then the consumer thread should handle it gracefully by briefly sleeping before polling again. The sleep is a small number, which you set to prevent busy waiting.\footnote{Busy waiting (or spinning) is when a consumer constantly asks an empty queue if it has any data in queue --- consequently, this tanks a systems performance.}
	
	You must wrap the procedure in a \code{multiple-value-bind} to check if the pop was successful.
	
\begin{lispcode}
(defparameter *work-queue* (make-queue))
;; Producer thread
(defun producer ()
(loop for i from 1 to 1000 do
(queue-push *work-queue* i)
(sleep 0.001)))

;; Consumer thread
(defun consumer ()
(loop
(multiple-value-bind (item found)
(queue-pop *work-queue*)
(if found
(format t "Processing: ~A~%" item)
(sleep 0.001)))))
\end{lispcode}
	
	\subsubsection{MPMC Bounded Queue (Vyukov)}
\begin{lispcode}
(defparameter *buffer* (make-bounded-queue 64))  ; Must be power of 2

;; Non-blocking operations
(when (bounded-queue-push *buffer* "data")
(format t "Push succeeded~%"))

(multiple-value-bind (value success)
(bounded-queue-pop *buffer*)
(when success
(format t "Got: ~A~%" value)))
\end{lispcode}
	
	\subsubsection{SPSC Bounded Queue}
\begin{lispcode}
(defparameter *spsc* (make-spsc-queue 256))  ; Must be power of 2
;; Producer thread only
(when (spsc-push *spsc* "fast-data")
(format t "Pushed successfully~%"))
;; Consumer thread only  
(multiple-value-bind (value success)
(spsc-pop *spsc*)
(when success
(format t "Got: ~A~%" value)))
\end{lispcode}
	
	\subsubsection{Error Handling}
	Operations return multiple values to indicate success/failure rather than signaling conditions:
\begin{lispcode}
;; Always check the second return value
(multiple-value-bind (value success)
(queue-pop *queue*)
(if success
(process value)
(handle-empty-queue)))
\end{lispcode}
	
	\section{API Reference}
	The API Reference for \texttt{cl-freelock} is a dictionary of all global functions, involving unbounded Multi-Producer Multi-Consumer (MPMC), bounded MPMC, bounded Single Producer Single Consumer (SPSC), error handling, and atomic primitives.\footnote{n.b, The difference between atomic primitives and atomic operations is a matter concerning layers of abstraction---operations are composed of combinations of primitives and non-primitive procedures.}
	
	\textbf{Additional Information}
	
	\begin{enumerate}
		\item \texttt{cl-freelock} allows the use of \code{fl} when calling its functions.
		\subitem  For example: \code{(fl:queue-push *work-queue* i)}
		\item The API reference for \texttt{cl-freelock} is mostly written with \textit{SBCL} in-mind. However, the differences between implementations regarding the code presented in this documentation will be quite minor.
	\end{enumerate}
	% ==========================
	\multicomment{According to ai: "This is the most important reason. When lwarp sees a description environment, it translates it directly into an HTML Definition List (<dl>, <dt>, <dd>). This is semantically perfect for the web. It means your API reference will automatically look great on mobile phones, and screen readers will perfectly understand the structure. (If you used Option 2 or 3, lwarp would generate an HTML <table>, which is terrible for mobile viewing)."
	I will need to investigate the HTML to confirm this is true, but for now it seems to make sense, so it wont hurt to atleast try. Additionally, it would stop the need for manual spacing with \\ and \vspace. I already import the enumitem package, so I might aswell extend its use.
	Update: This method DID NOT work, as lwarp can't handle enumitem's \item in certain contexts. Thus, we will just use standard latex instead.}
	% ==========================
	
	\subsection{Atomic Primitives}
	\subsubsection{ATOMIC-REF}
	\customitem{Signature:}{\code{atomic-ref}}
	
	\customitem{Description:}{A structure which holds a value that can be atomically updated--it is the foundational atomic pointer for \texttt{cl-freelock}.}

	\subsubsection{MAKE-ATOMIC-REF}
	
	\customitem{Signature:}{\code{(make-atomic-ref value)}}
	
	\customitem{Description:}{	Constructor for the \code{atomic-ref} structure---creates a new atomic reference that holds the provided \code{value}.}
	
	\subsubsection{ATOMIC-REF-VALUE}
	
	\customitem{Signature:}{\code{(atomic-ref-value atomic-ref)}}
	
	
	
	\customitem{Description:}{Accessor for the value held within an \code{atomic-ref}---used internally by Compare and Swap (CAS) operations.\footnote{An accessor (i.e., A structure or slot accessor) is a procedure which admits access to an object's internal state.}\supcom\footnote{n.b., \code{atomic-ref-value} is internally a \code{defstruct}. In \textit{Common Lisp}, the use of \code{defstruct} acts as both a getter and setter by the implementation's compiler. A getter is a procedure which specifically returns a read-only value.}}
	\subsubsection{CAS}
	
	\customitem{Signature:}{\code{(cas place old new)}}
	
	\customitem{Description:}{A CAS macro---atomically stored \code{new} in \code{place} if the current value of \code{place} is equal to \code{old}.}
	
	\customitem{Returns:}{The old value of \code{place}.}
	\subsubsection{ATOMIC-INCF}
	
	\customitem{Signature:}{\code{(atomic-incf atomic-ref \&optional (delta 1)}}
	
	\customitem{Description:}{Atomically increments the value of \code{atomic-ref} by \code{delta}. It is implemented via a CAS loop to facilitate compatibility across different \textit{Lisp} implementations.}
	
	\customitem{Returns:}{The new incremented value.}
	\subsubsection{MEMORY-BARRIER}
	
	\customitem{Signature:}{\code{(memory-barrier)}}
	
	\customitem{Description:}{Executes a memory barrier---ensures that all memory operations before the barrier are both completed and globally visible before any subsequent operations are allowed to begin.}
	\subsection{MPMC Unbounded Queue}
	
	\subsubsection{QUEUE}
	\customitem{Signature:}{\code{queue}}
	
	\customitem{Description:}{A lock-free, multi-producer, multi-consumer (MPMC) unbounded queue class based on the MichaelScott algorithm.}
	\subsubsection{MAKE-QUEUE}
	\customitem{Signature:}{\code{(make-queue)}}
	
	\customitem{Description:}{Creates an unbounded queue instance.}
	
	\customitem{Returns:}{A fresh \code{queue} object---initialized with a dummy node.}
	
	\customitem{Example:}{}
\begin{lispcode}
(defparameter *work-queue* (make-queue))
\end{lispcode}
	\subsubsection{QUEUE-PUSH}
	
	\customitem{Signature:}{}
	
	\customitem{Description:}{}
	\subsubsection{QUEUE-POP}
	
	\customitem{Signature:}{}
	
	\customitem{Description:}{}
	\subsubsection{QUEUE-EMPTY-P}
	
	\customitem{Signature:}{}
	
	\customitem{Description:}{}
	
	\subsection{MPMC Bounded Queue}
	\subsubsection{BOUNDED-QUEUE}
	
	\customitem{Signature:}{}
	
	\customitem{Description:}{}
	\subsubsection{MAKE-BOUNDED-QUEUE}
	
	\customitem{Signature:}{}
	
	\customitem{Description:}{}
	\subsubsection{BOUNDED-QUEUE-CAPACITY}
	
	\customitem{Signature:}{}
	
	\customitem{Description:}{}
	\subsubsection{BOUNDED-QUEUE-PUSH}
	
	\customitem{Signature:}{}
	
	\customitem{Description:}{}
	\subsubsection{BOUNDED-QUEUE-POP}
	
	\customitem{Signature:}{}
	
	\customitem{Description:}{}
	\subsubsection{BOUNDED-QUEUE-EMPTY-P}
	
	\customitem{Signature:}{}
	
	\customitem{Description:}{}
	\subsubsection{BOUNDED-QUEUE-FULL-P}
	
	\customitem{Signature:}{}
	
	\customitem{Description:}{}
	\subsubsection{BOUNDED-QUEUE-PUSH-BATCH}
	
	\customitem{Signature:}{}
	
	\customitem{Description:}{}
	\subsubsection{BOUNDED-QUEUE-POP-BATCH}
	
	\customitem{Signature:}{}
	
	\customitem{Description:}{}
	
	\subsection{SPSC Bounded Queue}
	\subsubsection{SPSC-QUEUE}
	
	\customitem{Signature:}{}
	
	\customitem{Description:}{}
	\subsubsection{MAKE-SPSC-QUEUE}
	
	\customitem{Signature:}{}
	
	\customitem{Description:}{}
	\subsubsection{SPSC-PUSH}
	
	\customitem{Signature:}{}
	
	\customitem{Description:}{}
	\subsubsection{SPSC-POP}
	
	\customitem{Signature:}{}
	
	\customitem{Description:}{}
	
	\section{Performance Guide}
	Tips and guidelines for optimal performance with \texttt{cl-freelock} --- considerations for the performance of unbounded Multi-Producer Multi-Consumer (MPMC), bounded MPMC, bounded Single Producer Single Consumer (SPSC), error handling, and atomic operations.
	
	\textbf{Additional Information}
	\begin{enumerate}
		\item \texttt{cl-freelock} allows the use of \code{fl} when calling its functions.
		\subitem  For example: \code{(fl:queue-push *work-queue* i)}
		\item The API reference for cl-freelock is mostly written with SBCL in-mind. However, the differences between implementations regarding the code presented in this documentation will be quite minor.
	\end{enumerate}
	\subsection{General Performance Principles}
	\subsubsection{Choose the Right Data Structure}
	\subsubsection{Capacity Selection}
	\subsubsection{Thread Contention}
	\subsection{Optimization Techniques}
	\subsubsection{Batch Operations}
	\subsubsection{Avoid Busy Waiting}
	\subsubsection{Memory Allocation}
	\subsection{Implementation-Specific Notes}
	\subsubsection{SBCL}
	\subsubsection{CLASP}
	\subsection{Common Pitfalls}
	\subsubsection{Wrong capacity}
	\subsubsection{Thread safety}
	\subsubsection{Busy waitings}
	\subsubsection{Memory leaks}
	\subsubsection{False sharing}
	\subsection{Recommended Patterns}
	\subsubsection{Producer-Consumer Pool}
	\subsubsection{Basic Producer-Consumer}
	\subsubsection{Bounded Buffer with Batch Processing}
	\subsubsection{SPSC Pipeline}
	\subsubsection{Load Balancing with Multiple Queues}
	\subsubsection{Web Server Request Queue}
	\subsubsection{Atomic Counter Service}
	\subsubsection{Event Processing Pipeline}
	\subsubsection{Memory Pool with Lock-Free Allocation}
	
	
\end{document}
	